\FigureLayoutDeclare{img/2007/beijing_science/q13_fig1.png}{5.0cm}{6bp 16bp 6bp 0bp}
\FigureLayoutDeclare{img/2007/beijing_science/q16_fig1.png}{6.0cm}{106.987bp 73.991bp 0bp 12.998bp}

\chapter{2007年北京卷（理）}
\section{单选题}
\begin{problem}
已知 \(\cos \theta \cdot \tan \theta < 0\)，那么角 \(\theta\) 是（\quad）

\choices
    {第一或第二象限角}
    {第二或第三象限角}
    {第三或第四象限角}
    {第一或第四象限角}
\end{problem}

\begin{answer}
C
\end{answer}

\begin{solution}
由于正切有定义，\(\cos\theta\ne0\)，所以
\(\cos\theta\tan\theta=\cos\theta\cdot\dfrac{\sin\theta}{\cos\theta}=\sin\theta\).
题设给出 \(\sin\theta<0\)，故 \(\theta\) 是第三或第四象限角，选 C.
\end{solution}

\begin{problem}
函数 \(f(x)=3^x\)（\(0 < x \leqslant 2\)）的反函数的定义域为（\quad）

\choices
    {\((0, +\infty)\)}
    {\((1, 9]\)}
    {\((0,1)\)}
    {\([9, +\infty)\)}
\end{problem}

\begin{answer}
B
\end{answer}

\begin{solution}
函数 \(f(x)=3^x\) 在 \(\R\) 上单调递增，因此当 \(0<x\leqslant2\) 时，
\(f(x)\) 的值域为 \((3^0,3^2]=(1,9]\). 反函数的定义域等于原函数的值域，故选 B.
\end{solution}

\begin{problem}
平面 \(\alpha \parallel \beta\) 的一个充分条件是（\quad）

\choices
    {存在一条直线 \(a\)，\(a \parallel \alpha\)，\(a \parallel \beta\)}
    {存在一条直线 \(a\)，\(a \subset \alpha\)，\(a \parallel \beta\)}
    {存在两条平行直线 \(a\)，\(b\)，\(a \subset \alpha\)，\(b \subset \beta\)，\(a \parallel \beta\)，\(b \parallel \alpha\)}
    {存在两条异面直线 \(a\)，\(b\)，\(a \subset \alpha\)，\(b \subset \beta\)，\(a \parallel \beta\)，\(b \parallel \alpha\)}
\end{problem}

\begin{answer}
D
\end{answer}

\begin{solution}
若两平面相交于直线 \(l\)，则平面 \(\alpha\) 内平行于 \(\beta\) 的直线只能平行于 \(l\)，平面 \(\beta\) 内平行于 \(\alpha\) 的直线也只能平行于 \(l\). 因而在这种情况下，两条直线必相互平行，不可能异面，所以选项 D 中的条件排除了两平面相交的情形，必有 \(\alpha\parallel\beta\).

选项 A 只说明存在同一条直线同时平行于两平面，另外两个非异面直线的选项也都不能排除两平面相交，故它们都不是充分条件.
\end{solution}

\begin{problem}
已知 \(O\) 是 \(\triangle ABC\) 所在平面内一点，\(D\) 为 \(BC\) 边中点，且 \(2\overrightarrow{OA} + \overrightarrow{OB} + \overrightarrow{OC} = 0\)，那么（\quad）

\choices
    {\(\overrightarrow{AO} = \overrightarrow{OD}\)}
    {\(\overrightarrow{AC} = 2\overrightarrow{OD}\)}
    {\(\overrightarrow{AC} = 3\overrightarrow{OD}\)}
    {\(2\overrightarrow{AO} = \overrightarrow{OD}\)}
\end{problem}

\begin{answer}
A
\end{answer}

\begin{solution}
因为 \(D\) 是 \(BC\) 的中点，所以
\(2\overrightarrow{OD}=\overrightarrow{OB}+\overrightarrow{OC}\).
由题设 \(2\overrightarrow{OA}+\overrightarrow{OB}+\overrightarrow{OC}=0\)，得
\(2\overrightarrow{OD}=-2\overrightarrow{OA}=2\overrightarrow{AO}\)，从而
\(\overrightarrow{OD}=\overrightarrow{AO}\)，故选 A.
\end{solution}

\begin{problem}
记者要为 \(5\) 名志愿者和他们帮助的 \(2\) 位老人拍照. 要求排成一排，\(2\) 位老人相邻但不得在两端. 不同的排法共有（\quad）

\choices
    {\(1440\) 种}
    {\(960\) 种}
    {\(720\) 种}
    {\(480\) 种}
\end{problem}

\begin{answer}
B
\end{answer}

\begin{solution}
把两位老人看作一个整体，则这个整体与 5 名志愿者共形成 6 个对象. 老人整体不能在两端，所以它在 6 个位置中有 \(4\) 个位置可选；两位老人内部有 \(2!\) 种排法，5 名志愿者有 \(5!\) 种排法. 因而排法总数为
\(4\cdot2!\cdot5!=960\)，故选 B.
\end{solution}

\begin{problem}
若不等式组 \( \begin{cases}
x-y&\geqslant0,\\
2x+y&\leqslant2,\\
y&\geqslant0,\\
x+y&\leqslant a
\end{cases} \) 表示的平面区域是一个三角形，则 \(a\) 的取值范围是（\quad）

\choices
    {\(a \geqslant \frac{4}{3}\)}
    {\(0 < a \leqslant 1\)}
    {\(1 \leqslant a \leqslant \frac{4}{3}\)}
    {\(0 < a \leqslant 1\) 或 \(a \geqslant \frac{4}{3}\)}
\end{problem}

\begin{answer}
D
\end{answer}

\begin{solution}
前三个不等式确定的三角形顶点为 \((0,0)\)、\((1,0)\)、\(\left(\dfrac23,\dfrac23\right)\). 新边界为 \(x+y=a\).

当 \(0<a\leqslant1\) 时，新边界分别与 \(y=x\)、\(y=0\) 相交，所得区域仍为三角形；当 \(1<a<\dfrac43\) 时，新边界截去原三角形的一个顶点，所得区域为四边形；当 \(a\geqslant\dfrac43\) 时，新不等式不再截去三角形的内部，所得区域仍为原三角形. 当 \(a\leqslant0\) 时区域至多退化，不能成为三角形.

所以 \(a\) 的取值范围为 \(0<a\leqslant1\) 或 \(a\geqslant\dfrac43\)，故选 D.
\end{solution}

\begin{problem}
如果正数 \(a\)，\(b\)，\(c\)，\(d\) 满足 \(a + b = cd = 4\). 那么（\quad）

\choices
    {\(ab \leqslant c + d\)，且等号成立时 \(a, b, c, d\) 的取值唯一}
    {\(ab \geq c + d\)，且等号成立时 \(a\)，\(b\)，\(c\)，\(d\) 的取值唯一}
    {\(ab \leqslant c + d\)，且等号成立时 \(a\)，\(b\)，\(c\)，\(d\) 的取值不唯一}
    {\(ab \geqslant c + d\)，且等号成立时 \(a\)，\(b\)，\(c\)，\(d\) 的取值不唯一}
\end{problem}

\begin{answer}
A
\end{answer}

\begin{solution}
由基本不等式，
\(ab\leqslant\left(\dfrac{a+b}{2}\right)^2=4\)，且等号当且仅当 \(a=b=2\).
又因为 \(c\)，\(d>0\) 且 \(cd=4\)，所以
\(c+d\geqslant2\sqrt{cd}=4\)，且等号当且仅当 \(c=d=2\).
于是 \(ab\leqslant4\leqslant c+d\). 两个不等式同时取等号时必须有 \(a=b=c=d=2\)，取值唯一，故选 A.
\end{solution}

\begin{problem}
对于函数
\begin{circlelist}
\item \(f(x)=\lg(|x-2|+1)\)；
\item \(f(x)=(x-2)^2\)；
\item \(f(x)=\cos(x+2)\).
\end{circlelist}
判断如下三个命题的真假：

命题甲：\(f(x+2)\) 是偶函数；

命题乙：\(f(x)\) 在 \((-\infty,2)\) 上是减函数，在 \((2,+\infty)\) 上是增函数；

命题丙：\(f(x+2)-f(x)\) 在 \((-\infty,+\infty)\) 上是增函数. 

能使命题甲，乙，丙均为真的所有函数的序号是（\quad）

\choices
    {\circled{1}\circled{3}}
    {\circled{1}\circled{2}}
    {\circled{3}}
    {\circled{2}}
\end{problem}

\begin{answer}
D
\end{answer}

\begin{solution}
对第 1 个函数，\(f(x+2)=\lg(|x|+1)\) 是偶函数，且原函数在 \((-\infty,2)\) 上递减、在 \((2,+\infty)\) 上递增；但是当 \(x<0\) 时，
\[
f(x+2)-f(x)=\lg\dfrac{1-x}{3-x}
\]
该式随 \(x\) 递减，所以命题丙不真.

对第 2 个函数，\(f(x+2)=x^2\) 是偶函数；\(f(x)=(x-2)^2\) 在 \((-\infty,2)\) 上递减、在 \((2,+\infty)\) 上递增；并且
\[
f(x+2)-f(x)=x^2-(x-2)^2=4x-4
\]
这是增函数，所以三个命题均真.

对第 3 个函数，\(f(x+2)=\cos(x+4)\) 不是偶函数，且余弦函数也不在所给两个区间上分别保持单调，故三个命题不可能均真. 因而只有序号 2，选 D.
\end{solution}

\section{填空题}
\begin{problem}
\(\frac{2}{(1 + \i)^2} =\)\fillinblank{}.
\end{problem}

\begin{answer}
\(-\i\)
\end{answer}

\begin{solution}
\((1+\i)^2=1+2\i+\i^2=2\i\)，所以
\[
\frac{2}{(1+\i)^2}=\frac{2}{2\i}=\frac1\i=-\i
\]
\end{solution}

\begin{problem}
若数列 \( \{a_{n}\} \) 的前 \(n\) 项和 \(S_{n}=n^{2}-10n\)（\(n=1\)，\(2\)，\(3\)，\(\cdots\)），则此数列的通项公式为\fillinblank{}；数列 \( \{na_{n}\} \) 中数值最小的项是第\fillinblank{} 项.
\end{problem}

\begin{answer}
\(2n-11\)，\(3\)
\end{answer}

\begin{solution}
由 \(a_1=S_1=1-10=-9\)，且当 \(n\geqslant2\) 时
\[
a_n=S_n-S_{n-1}=n^2-10n-\bigl((n-1)^2-10(n-1)\bigr)=2n-11
\]
当 \(n=1\) 时该式也给出 \(-9\)，故 \(a_n=2n-11\).

令 \(b_n=na_n=2n^2-11n\)，则
\(b_{n+1}-b_n=4n-9\). 该差在 \(n=1\)，\(2\) 时为负，在 \(n\geqslant3\) 时为正，所以 \(b_n\) 先减后增，最小项为 \(b_3\)，即第 3 项.
\end{solution}

\begin{problem}
在 \(\triangle ABC\) 中，若 \(\tan A = \frac{1}{3}\)，\(C = 150^{\circ}\)，\(BC = 1\)，则 \(AB =\fillinblank{}\).
\end{problem}

\begin{answer}
\(\dfrac{\sqrt{10}}{2}\)
\end{answer}

\begin{solution}
三角形内角和给出 \(A+B=30^\circ\). 因为 \(\tan A=\dfrac13>0\) 且 \(A<30^\circ\)，所以
\(\sin A=\dfrac1{\sqrt{10}}\).
由正弦定理，
\(\frac{AB}{\sin C}=\frac{BC}{\sin A}\)，\(AB=\frac{\sin150^\circ}{\sin A} =\frac{1/2}{1/\sqrt{10}}=\frac{\sqrt{10}}2\).
\end{solution}

\begin{problem}
已知集合 \(A = \{x \mid |x - a| \leqslant 1\}\)，\(B = \{x \mid x^2 - 5x + 4 \geqslant 0\}\). 若 \(A \cap B = \varnothing\)，则实数 \(a\) 的取值范围是\fillinblank{}.
\end{problem}

\begin{answer}
\((2,3)\)
\end{answer}

\begin{solution}
\(x^2-5x+4=(x-1)(x-4)\)，所以
\(B=(-\infty,1]\cup[4,+\infty)\). 又 \(A=[a-1,a+1]\).
要使 \(A\cap B=\varnothing\)，闭区间 \([a-1,a+1]\) 必须完全位于开区间 \((1,4)\) 内，因此
\(a-1>1\)，且 \(a+1<4\). 从而 \(2<a<3\).
\end{solution}

\begin{problem}
2002 年在北京召开的国际数学家大会，会标是以我国古代数学家赵爽的弦图为基础设计的. 弦图是由四个全等直角三角形与一个小正方形拼成的一个大正方形（如图）.

\bitmapfigure[max width=0.24\linewidth]{img/2007/beijing_science/q13_fig1.png}

如果小正方形的面积为 \(1\)，大正方形的面积为 \(25\)，直角三角形中较小的锐角为 \(\theta\)，那么 \(\cos 2\theta\) 的值等于\fillinblank{}.
\end{problem}

\begin{answer}
\(\dfrac{7}{25}\)
\end{answer}

\begin{solution}
设弦图中直角三角形的斜边为大正方形的边，长度为 \(5\)，两直角边分别为 \(5\cos\theta\) 和 \(5\sin\theta\). 因为 \(\theta\) 是较小的锐角，所以 \(0<\theta<\dfrac\pi4\)，且小正方形的边长为
\(5\cos\theta-5\sin\theta=1\).
于是 \(\cos\theta-\sin\theta=\dfrac15\)，两边平方得
\(1-\sin2\theta=\frac1{25}\)，\(\sin2\theta=\frac{24}{25}\).
又 \(0<2\theta<\dfrac\pi2\)，故 \(\cos2\theta>0\)，从而
\[
\cos2\theta=\sqrt{1-\sin^22\theta}
=\sqrt{1-\frac{576}{625}}=\frac7{25}
\]
\end{solution}

\begin{problem}
已知函数 \(f(x)\)，\(g(x)\) 分别由下表给出

\begin{center}
\begin{tabular}{c|ccc}
\(x\) & \(1\) & \(2\) & \(3\) \\
\hline
\(f(x)\) & \(1\) & \(3\) & \(1\)
\end{tabular}
\hspace{1.5em}
\begin{tabular}{c|ccc}
\(x\) & \(1\) & \(2\) & \(3\) \\
\hline
\(g(x)\) & \(3\) & \(2\) & \(1\)
\end{tabular}
\end{center}

则 \(f[g(1)]\) 的值为\fillinblank{}；满足 \(f[g(x)] > g[f(x)]\) 的 \(x\) 的值是\fillinblank{}.
\end{problem}

\begin{answer}
\(1\)，\(2\)
\end{answer}

\begin{solution}
由表直接计算，\(g(1)=3\)，所以 \(f[g(1)]=f(3)=1\).

逐一检验第二个条件：当 \(x=1\) 时，\(f[g(1)]=1<3=g[f(1)]\)；当 \(x=2\) 时，\(f[g(2)]=f(2)=3>1=g[f(2)]\)；当 \(x=3\) 时，\(f[g(3)]=f(1)=1<3=g[f(3)]\). 因而满足条件的 \(x\) 只有 \(2\).
\end{solution}

\section{解答题}
\begin{problem}
数列 \( \{a_{n}\} \) 中，\(a_{1}=2\)，\(a_{n+1}=a_{n}+cn\)（\(c\) 是常数，\(n=1\)，\(2\)，\(3\)，\(\cdots\)）. 且 \(a_{1}\)，\(a_{2}\)，\(a_{3}\) 成公比不为 \(1\) 的等比数列.

\begin{enumerate}
\item 求 \(c\) 的值.
\item 求 \( \{a_{n}\} \) 的通项公式.
\end{enumerate}
\end{problem}

\begin{solution}
\begin{enumerate}
\item 由递推式，\(a_2=2+c\)，\(a_3=a_2+2c=2+3c\). 因为前三项成等比数列且公比不为 \(1\)，有
\[
(2+c)^2=a_1a_3=2(2+3c)
\]
整理得 \(c(c-2)=0\). 若 \(c=0\)，则前三项及以后各项均为 \(2\)，公比为 \(1\)，不合题意，所以 \(c=2\).
\item 当 \(c=2\) 时，
\[
a_n=a_1+\sum_{j=1}^{n-1}2j
=2+n(n-1)=n^2-n+2
\]
\end{enumerate}
\end{solution}

\begin{problem}
如图，在 \(\mathrm{Rt}\triangle AOB\) 中，\(\angle OAB = \frac{\pi}{6}\)，斜边 \(AB = 4\). \(\mathrm{Rt}\triangle AOC\) 可以通过 \(\mathrm{Rt}\triangle AOB\) 以直线 \(AO\) 为轴旋转得到，且二面角 \(B-AO-C\) 是直二面角. 动点 \(D\) 在斜边 \(AB\) 上.

\begin{enumerate}
\item 求证：平面 \(COD\) \(\perp\) 平面 \(AOB\).
\item 当 \(D\) 为 \(AB\) 的中点时，求异面直线 \(AO\) 与 \(CD\) 所成角的大小.
\item 求 \(CD\) 与平面 \(AOB\) 所成角的最大值.
\end{enumerate}

\bitmapfigure[max width=0.34\linewidth]{img/2007/beijing_science/q16_fig1.png}
\end{problem}

\begin{solution}
\begin{enumerate}
\item 在直角三角形 \(AOB\) 中，
\(AO=AB\cos\frac\pi6=2\sqrt3\)，\(OB=AB\sin\frac\pi6=2\).
旋转关系给出 \(OC=OB=2\) 且 \(OC\perp AO\). 又因为二面角 \(B-AO-C\) 是直二面角，所以 \(OB\perp OC\). 直线 \(OC\) 垂直于平面 \(AOB\) 内相交的两条直线 \(AO\)、\(OB\)，故 \(OC\perp\) 平面 \(AOB\). 由于 \(OC\subset\) 平面 \(COD\)，所以平面 \(COD\perp\) 平面 \(AOB\).
\item 当 \(D\) 为 \(AB\) 中点时，在平面 \(AOB\) 内过 \(D\) 作 \(DE\perp OB\)，交 \(OB\) 于 \(E\). 因为 \(AO\perp OB\)，所以 \(DE\parallel AO\)，从而 \(E\) 是 \(OB\) 的中点，
\(OE=1\)，\(DE=\frac{AO}{2}=\sqrt3\).
由 \(OC\perp\) 平面 \(AOB\) 且 \(OC=2\)，得
\(CE=\sqrt{OC^2+OE^2}=\sqrt5\). 又 \(CE\perp DE\)，所以异面直线 \(AO\) 与 \(CD\) 所成的角等于 \(\angle CDE\)，且
\[
\tan\angle CDE=\frac{CE}{DE}=\sqrt{\frac53}=\frac{\sqrt{15}}3
\]
故所求角为 \(\arctan\dfrac{\sqrt{15}}3\).
\item 设 \(H\) 是点 \(O\) 到斜边 \(AB\) 的垂足. 对任意 \(D\in AB\)，直线 \(CD\) 在平面 \(AOB\) 上的射影为 \(OD\)，所以若其与平面 \(AOB\) 所成角为 \(\alpha\)，则
\[
\tan\alpha=\frac{OC}{OD}=\frac2{OD}
\]
当 \(D=H\) 时，\(OD\) 最小. 由直角三角形面积或射影定理，
\[
OH=\frac{AO\cdot OB}{AB}=\frac{2\sqrt3\cdot2}{4}=\sqrt3
\]
因此 \(\tan\alpha\) 的最大值为 \(\dfrac2{\sqrt3}=\dfrac{2\sqrt3}3\)，所求最大角为 \(\arctan\dfrac{2\sqrt3}3\).
\end{enumerate}
\end{solution}

\begin{problem}
矩形 \(ABCD\) 的两条对角线相交于点 \(M(2,0)\)，\(AB\) 边所在直线的方程为 \(x - 3y - 6 = 0\)，点 \(T(-1,1)\) 在 \(AD\) 边所在直线上.

\begin{enumerate}
\item 求 \(AD\) 边所在直线的方程.
\item 求矩形 \(ABCD\) 外接圆的方程.
\item 若动圆 \(P\) 过点 \(N(-2,0)\)，且与矩形 \(ABCD\) 的外接圆外切，求动圆 \(P\) 的圆心的轨迹方程.
\end{enumerate}
\end{problem}

\begin{solution}
\begin{enumerate}
\item 直线 \(AB\) 的斜率为 \(\dfrac13\)，而矩形相邻边垂直，所以直线 \(AD\) 的斜率为 \(-3\). 过 \(T(-1,1)\) 得
\(y-1=-3(x+1)\)，\(3x+y+2=0\).
\item 点 \(A\) 是直线 \(AB\) 与 \(AD\) 的交点. 联立
\(x-3y-6=0\) 与 \(3x+y+2=0\)，得 \(A(0,-2)\). 矩形两条对角线交点 \(M(2,0)\) 是外接圆圆心，且
\[
MA^2=(0-2)^2+(-2)^2=8
\]
所以外接圆方程为 \((x-2)^2+y^2=8\).
\item 设动圆圆心为 \(P(x,y)\). 动圆过 \(N(-2,0)\)，半径为 \(PN\)；与圆心 \(M(2,0)\)、半径 \(2\sqrt2\) 的外接圆外切，所以
\(PM=PN+2\sqrt2\)，\(PM-PN=2\sqrt2\).
这是以 \(M\)，\(N\) 为焦点、焦距差为 \(2\sqrt2\) 的双曲线. 其焦点距为 \(4\)，故半长轴为 \(\sqrt2\)，半短轴满足 \(b^2=4-2=2\)，轨迹所在双曲线为
\[
\frac{x^2}{2}-\frac{y^2}{2}=1
\]
由于 \(PM>PN\)，点 \(P\) 在左支，故还需满足 \(x\leqslant-\sqrt2\).
\end{enumerate}
\end{solution}

\begin{problem}
某中学号召学生在今年春节期间至少参加一次社会公益活动（以下简称活动）. 该校合唱团共有 \(100\) 名学生，他们参加活动的次数统计如图所示.

\begin{enumerate}
\item 求合唱团学生参加活动的人均次数.
\item 从合唱团中任意选两名学生，求他们参加活动次数恰好相等的概率.
\item 从合唱团中任选两名学生，用 \(\xi\) 表示选两人参加活动次数之差的绝对值，求随机变量 \(\xi\) 的分布列及数学期望 \(E(\xi)\).
\end{enumerate}

\bitmapfigure[max width=0.38\linewidth]{img/2007/beijing_science/q18_fig1.png}
\end{problem}

\begin{solution}
\begin{enumerate}
\item 由图可知参加 1 次、2 次、3 次活动的学生人数分别为 \(10\)，\(50\)，\(40\)，所以人均次数为
\[
\frac{1\cdot10+2\cdot50+3\cdot40}{100}=\frac{230}{100}=2.3
\]
\item 从 100 名学生中任意选 2 名，样本总数为 \(\mathrm C_{100}^{2}=4950\). 两人的次数相等的选法数为
\(\mathrm C_{10}^{2}+\mathrm C_{50}^{2}+\mathrm C_{40}^{2}=45+1225+780=2050\)，故所求概率为
\[
\frac{2050}{4950}=\frac{41}{99}
\]
\item 两人次数之差的绝对值只能取 \(0\)，\(1\)，\(2\). 于是
\[
P(\xi=0)=\frac{\mathrm C_{10}^{2}+\mathrm C_{50}^{2}+\mathrm C_{40}^{2}}{\mathrm C_{100}^{2}}=\frac{41}{99}
\]
\(P(\xi=1)=\frac{10\cdot50+50\cdot40}{\mathrm C_{100}^{2}}=\frac{50}{99}\)，\(P(\xi=2)=\frac{10\cdot40}{\mathrm C_{100}^{2}}=\frac8{99}\).
分布列为
\[
\begin{tabular}{c|ccc}
\(\xi\)&\(0\)&\(1\)&\(2\)\\ \hline
\(P\)&\(\dfrac{41}{99}\)&\(\dfrac{50}{99}\)&\(\dfrac8{99}\)
\end{tabular}
\]
数学期望为
\[
E(\xi)=0\cdot\frac{41}{99}+1\cdot\frac{50}{99}+2\cdot\frac8{99}=\frac{66}{99}=\frac23
\]
\end{enumerate}
\end{solution}

\begin{problem}
如图，有一块半椭圆形钢板，其长半轴长为 \(2r\)，短半轴长为 \(r\). 计划将此钢板切割成等腰梯形的形状，下底 \(AB\) 是半椭圆的短轴，上底 \(CD\) 的端点在椭圆上，记 \(CD = 2x\)，梯形面积为 \(S\).

\begin{enumerate}
\item 求面积 \(S\) 以 \(x\) 为自变量的函数式，并写出其定义域.
\item 求面积 \(S\) 的最大值.
\end{enumerate}

\bitmapfigure[max width=0.34\linewidth]{img/2007/beijing_science/q19_fig1.png}
\end{problem}

\begin{solution}
\begin{enumerate}
\item 取半椭圆短轴所在直线为 \(x\) 轴，椭圆中心为原点，且取上半椭圆，则椭圆方程为
\(\dfrac{x^2}{r^2}+\dfrac{y^2}{4r^2}=1\). 由于 \(CD=2x\)，其两个端点的横坐标为 \(-x\)、\(x\)，纵坐标均为 \(2\sqrt{r^2-x^2}\)，而 \(AB=2r\). 所以梯形高为 \(2\sqrt{r^2-x^2}\)，面积为
\[
S=\frac{(2r+2x)\sqrt{r^2-x^2}}2=2(r+x)\sqrt{r^2-x^2}
\]
为得到非退化梯形，\(0<x<r\)，这就是定义域.
\item 在定义域内求导，
\[
S'(x)=2\frac{r^2-rx-2x^2}{\sqrt{r^2-x^2}}
=2\frac{(r-2x)(r+x)}{\sqrt{r^2-x^2}}
\]
因为 \(r+x>0\)，所以当 \(0<x<\dfrac r2\) 时 \(S'(x)>0\)，当 \(\dfrac r2<x<r\) 时 \(S'(x)<0\). 因而 \(x=\dfrac r2\) 时面积取得最大值，
\[
S_{\max}=2\left(r+\frac r2\right)\sqrt{r^2-\frac{r^2}{4}}
=\frac{3\sqrt3}{2}r^2
\]
\end{enumerate}
\end{solution}

\begin{problem}
已知集合 \(A = \{a_{1}, a_{2}, \cdots, a_{k}\}\)（\(k \geqslant 2\)），其中 \(a_{i} \in \Z\)（\(i = 1\)，\(2\)，\(\cdots\)，\(k\)）. 由 \(A\) 中的元素构成两个相应的集合：\(S = \{(a, b) \mid a \in A, b \in A, a + b \in A\}\)；\(T = \{(a, b) \mid a \in A, b \in A, a - b \in A\}\)，其中 \((a, b)\) 是有序数对，集合 \(S\) 和 \(T\) 中的元素个数分别为 \(m\) 和 \(n\). 若对于任意的 \(a \in A\)，总有 \(-a \notin A\)，则称集合 \(A\) 具有性质 \(P\).

\begin{enumerate}
\item 检验集合 \(\{0, 1, 2, 3\}\) 与 \(\{-1, 2, 3\}\) 是否具有性质 \(P\)，并对其中具有性质 \(P\) 的集合，写出相应的集合 \(S\) 和 \(T\).
\item 对任何具有性质 \(P\) 的集合 \(A\)，证明：\(n \leqslant \frac{k(k - 1)}{2}\).
\item 判断 \(m\) 和 \(n\) 的大小关系，并证明你的结论.
\end{enumerate}
\end{problem}

\begin{solution}
\begin{enumerate}
\item 集合 \(\{0,1,2,3\}\) 中有 \(0\)，而 \(-0=0\) 仍在集合中，所以它不具有性质 \(P\). 对集合 \(A=\{-1,2,3\}\)，\(-(-1)=1\)、\(-2=-2\)、\(-3=-3\) 均不在 \(A\) 中，因此它具有性质 \(P\).

对该集合逐项检验和与差，得到
\(S=\{(-1,3),(3,-1)\}\)，\(T=\{(2,-1),(2,3)\}\).
\item 对任意 \((a,b)\in T\)，有 \(a-b\in A\). 若 \(a=b\)，则 \(0\in A\)，这与性质 \(P\) 矛盾，所以 \(a\ne b\). 又不可能同时有 \((a,b)\in T\) 和 \((b,a)\in T\)，因为这将导致 \(a-b\) 与 \(b-a=-(a-b)\) 同时属于 \(A\)，也与性质 \(P\) 矛盾.

因此每一对不同元素至多贡献一个有序数对，故
\[
n\leqslant\mathrm C_k^2=\frac{k(k-1)}2
\]
\item 定义映射 \(\varphi:S\to T\)，\(\varphi(a,b)=(a+b,b)\). 若 \((a,b)\in S\)，则 \(a+b\in A\)、\(b\in A\)，且 \((a+b)-b=a\in A\)，所以 \(\varphi(a,b)\in T\).

反之，定义 \(\psi:T\to S\)，\(\psi(c,d)=(c-d,d)\). 若 \((c,d)\in T\)，则 \(c-d\in A\)、\(d\in A\)，并且 \((c-d)+d=c\in A\)，所以 \(\psi(c,d)\in S\). 显然 \(\psi\) 与 \(\varphi\) 互为逆映射，因此 \(S\) 与 \(T\) 等势，得到 \(m=n\).
\end{enumerate}
\end{solution}
